\documentclass{beamer}

% Theme choice
\usetheme{Madrid} % You can change to e.g., Warsaw, Berlin, CambridgeUS, etc.

% Encoding and font
\usepackage[utf8]{inputenc}
\usepackage[T1]{fontenc}

% Graphics and tables
\usepackage{graphicx}
\usepackage{booktabs}

% Code listings
\usepackage{listings}
\lstset{
  language=Python,
  basicstyle=\ttfamily\small,
  keywordstyle=\color{blue},
  commentstyle=\color{gray}\itshape,
  stringstyle=\color{red},
  breaklines=true,
  frame=single,
  showstringspaces=false
}

% Math packages
\usepackage{amsmath}
\usepackage{amssymb}

% Colors
\usepackage{xcolor}

% TikZ and PGFPlots
\usepackage{tikz}
\usepackage{pgfplots}
\pgfplotsset{compat=1.18}
\usetikzlibrary{positioning}

% Hyperlinks
\usepackage{hyperref}

% Title information
\title{Chapter 4: Data Structures I: Lists}
\author{Teaching Assistant}
\institute{University Name}
\date{\today}

\begin{document}

\frame{\titlepage}

\begin{frame}[fragile]
    \frametitle{Chapter 4: Data Structures I -- The Problem}

    \begin{block}{The Challenge: Managing Multiple Related Values}
        So far, we've used variables to store single pieces of information:
        \begin{lstlisting}[language=Python, basicstyle=\small\ttfamily]
student_name = "Alice"
final_score = 95
        \end{lstlisting}
        But what if we need to handle a collection of related data, like scores for an entire class?
    \end{block}
    
    \pause

    \begin{alertblock}{An Inefficient Approach}
        Creating a separate variable for each value is problematic:
        \begin{lstlisting}[language=Python, basicstyle=\small\ttfamily]
score1 = 88
score2 = 92
score3 = 77
# ...what if we had 30 students? Or 100?
        \end{lstlisting}
    \end{alertblock}

    \pause
    
    \begin{itemize}
        \item \textbf{Inconvenient:} Tedious to create and manage many variables.
        \item \textbf{Not Scalable:} Adding new data requires changing the code (adding new variables).
        \item \textbf{Hard to Process:} How would you calculate the average? You'd need to name every single variable.
    \end{itemize}

\end{frame}

\begin{frame}[fragile]
    \frametitle{Chapter 4: Data Structures I -- The Solution: Lists}

    \begin{block}{A Better Way: Using a Data Structure}
        We need a way to group a collection of items into a \textbf{single variable}. This is the purpose of a \textbf{data structure}.
        \medskip
        The first and most common data structure we'll learn is the \textbf{List}.
    \end{block}

    \pause

    \begin{block}{Analogy: From a Single Box to a Shelf}
        \begin{itemize}
            \item \textbf{A Variable} is like a single, labeled box.
            \begin{itemize}
                \item It has one name (e.g., \texttt{score1}).
                \item It holds one value (e.g., \texttt{88}).
            \end{itemize}
            \medskip
            \item \textbf{A List} is like a shelf with many numbered compartments.
            \begin{itemize}
                \item The entire shelf has one name (e.g., \texttt{scores}).
                \item It holds an entire collection of values, each in its own spot.
            \end{itemize}
        \end{itemize}
    \end{block}
    
    \pause
    
    \begin{center}
        \Large \texttt{scores = [88, 92, 77, 95, 100]}
    \end{center}

    \vfill
    This single \texttt{scores} variable gives us a powerful and organized way to manage our collection of data.

\end{frame}

\begin{frame}[fragile]
    \frametitle{What is a List? - Core Properties}
    
    A list is a container that holds a collection of items in a specific sequence. It is defined by three key characteristics:
    
    \vspace{1em}
    
    \begin{description}
        \item[\textbf{Ordered}] 
        Items have a defined sequence that doesn't change on its own. The first item you add stays the first item.
        \begin{itemize}
            \item \textit{Analogy: Like a numbered to-do list or people standing in a line.}
        \end{itemize}
        
        \item[\textbf{Mutable}] 
        "Mutable" means changeable. You can \textbf{add}, \textbf{remove}, or \textbf{change} items after the list has been created.
        
        \item[\textbf{Collection}] 
        A list can hold multiple items in a single variable. These items can be of different data types (e.g., numbers, strings, etc.).
    \end{description}
    
    \begin{alertblock}{Key Idea}
        A list is your go-to tool for storing a flexible, ordered sequence of items.
    \end{alertblock}
    
\end{frame}

\begin{frame}[fragile]
    \frametitle{What is a List? - Syntax \& Examples}
    
    \textbf{Syntax}: A list is created with square brackets \texttt{[]}, with items inside separated by commas.
    
    \vspace{1em}
    
    \begin{block}{Examples}
\begin{lstlisting}[language=Python]
# A list of integers
student_scores = [88, 92, 77, 95, 100]

# A list of strings
shopping_list = ["apples", "milk", "bread"]

# A list can hold mixed data types
# (string, integer, float, boolean)
employee_data = ["John Doe", 45, 85500.75, True]

# You can start with an empty list
empty_list = []
\end{lstlisting}
    \end{block}
    
    \vfill
    
    \textbf{Takeaway}: Remember these three words: \textbf{Ordered, Mutable, Collection}. On the next slide, we will see how the "ordered" property lets us access any item we want.
    
\end{frame}

\begin{frame}[fragile]
    \frametitle{Accessing Elements: The Concept of Indexing}
    
    Since lists are an \textbf{ordered} collection, every item has a numbered position, called its \textbf{index}. We use the index to look up a specific item.
    
    \vspace{1em}
    
    \begin{alertblock}{The Golden Rule: Counting Starts at Zero}
        Python lists are \textbf{zero-indexed}. This is a crucial concept.
        \begin{itemize}
            \item The \textbf{first} item is at index \texttt{0}.
            \item The \textbf{second} item is at index \texttt{1}.
            \item The \textbf{third} item is at index \texttt{2}, and so on.
        \end{itemize}
    \end{alertblock}
    
    \vspace{1em}
    
    \begin{block}{A Helpful Mental Model}
        The index of an item is the number of items that come \textit{before} it.
    \end{block}
    
\end{frame}

\begin{frame}[fragile]
    \frametitle{Accessing Elements: Syntax \& Example}

    \begin{block}{Access Syntax}
        Use the list's name followed by the index in square brackets \texttt{[]}.
        \begin{center}
            \texttt{list\_name[index]}
        \end{center}
    \end{block}

    \pause

    \begin{exampleblock}{Example Walkthrough}
\begin{lstlisting}[language=Python]
# 1. Define a list of student scores
student_scores = [88, 92, 77, 95]

# 2. Get the *first* score (at index 0)
first_score = student_scores[0]
# first_score is now 88

# 3. Get the *third* score (at index 2)
third_score = student_scores[2]
# third_score is now 77
\end{lstlisting}
    \end{exampleblock}

\end{frame}

\begin{frame}[fragile]
    \frametitle{Accessing Elements: Visualization \& Common Pitfalls}
    
    \begin{block}{Visualizing Indices for \texttt{student\_scores}}
        \centering
        \renewcommand{\arraystretch}{1.5}
        \begin{tabular}{|l|c|c|c|c|}
            \hline
            \textbf{Value} & \texttt{88} & \texttt{92} & \texttt{77} & \texttt{95} \\
            \hline
            \textbf{Index} & \texttt{0} & \texttt{1} & \texttt{2} & \texttt{3} \\
            \hline
        \end{tabular}
    \end{block}
    
    \pause
    
    \begin{alertblock}{Common Pitfall: The "Off-By-One" Error}
        New programmers often forget to start at zero.
        \begin{itemize}
            \item Trying to get the first item with \texttt{student\_scores[1]} will actually retrieve the \textbf{second} item (92).
            \item \textbf{Always remember: First item = index 0.}
        \end{itemize}
    \end{alertblock}
    
\end{frame}

\begin{frame}[fragile]
    \frametitle{Visualizing Indices}
    
    Let's visualize our list, \texttt{student\_scores = [88, 92, 77, 95]}.
    \vspace{0.5cm}
    
    Each value has a corresponding index that we use to access it.
    
    \begin{center}
    \begin{tabular}{|c|c|c|c|c|}
        \hline
        \textbf{Value} & 88 & 92 & 77 & 95 \\
        \hline
        \textbf{Index} & 0 & 1 & 2 & 3 \\
        \hline
    \end{tabular}
    \end{center}
    
    \pause
    \vspace{0.5cm}
    
    \begin{alertblock}{IndexError: Out of Bounds}
        The valid indices for this list are 0, 1, 2, and 3. Trying to access any other index will result in an \texttt{IndexError}.
    \end{alertblock}
    
    \begin{exampleblock}{Code Example: Causing an Error}
\begin{lstlisting}[language=Python]
# This list has 4 items, so the last index is 3.
# Trying to access index 4 will fail.
print(student_scores[4])

# Output: IndexError: list index out of range
\end{lstlisting}
    \end{exampleblock}
    
\end{frame}

\begin{frame}[fragile]
    \frametitle{Visualizing Indices - Negative Indexing}
    
    Python offers a convenient shortcut: \textbf{negative indexing} to access items from the \textit{end} of the list.
    
    \begin{itemize}
        \item \texttt{[-1]} refers to the \textbf{last} item.
        \item \texttt{[-2]} refers to the \textbf{second-to-last} item, and so on.
    \end{itemize}
    
    \pause
    
    Let's update our visualization to include negative indices:
    \begin{center}
    \begin{tabular}{|c|c|c|c|c|}
        \hline
        \textbf{Value} & 88 & 92 & 77 & 95 \\
        \hline
        \textbf{Positive Index} & 0 & 1 & 2 & 3 \\
        \hline
        \textbf{Negative Index} & -4 & -3 & -2 & -1 \\
        \hline
    \end{tabular}
    \end{center}
    
    \begin{exampleblock}{Code Example: Using Negative Indices}
\begin{lstlisting}[language=Python]
student_scores = [88, 92, 77, 95]

# Access the last score
last_score = student_scores[-1]
print(f"Last score: {last_score}") # Output: Last score: 95

# Access the second-to-last score
second_last = student_scores[-2]
print(f"Second to last: {second_last}") # Output: Second to last: 77
\end{lstlisting}
    \end{exampleblock}
    
\end{frame}

\begin{frame}[fragile]
    \frametitle{Modifying List Elements: The Concept}
    
    One of the most important features of a Python list is that it is \textbf{mutable}.
    
    \begin{itemize}
        \item \textbf{Mutable}: The object can be changed \textit{after} it has been created.
        \item This means you can change individual elements directly without creating a new list.
    \end{itemize}
    
    \vfill
    
    \begin{block}{Syntax for Modification}
        You change an item by accessing its index and assigning a new value.
        \begin{center}
            \texttt{list\_name[index] = new\_value}
        \end{center}
    \end{block}
    
\end{frame}

\begin{frame}[fragile]
    \frametitle{Modifying List Elements: An Example}
    
    Let's correct an item on our shopping list.
    
    \textbf{1. Initial List:}
    \begin{lstlisting}[language=Python]
shopping_list = ['milk', 'eggs', 'bread']
    \end{lstlisting}
    
    \begin{center}
    \begin{tabular}{|c|c|c|c|}
        \hline
        \textbf{Value} & \texttt{'milk'} & \texttt{'eggs'} & \texttt{'bread'} \\
        \hline
        \textbf{Index} & 0 & 1 & 2 \\
        \hline
    \end{tabular}
    \end{center}
    
    \pause
    
    \vfill
    
    \textbf{2. Perform Modification:}
    \begin{lstlisting}[language=Python]
# Oops, we need whole wheat bread
shopping_list[2] = 'whole wheat bread'
    \end{lstlisting}
    
    \pause
    
    \textbf{3. Result (Modified In-Place):}
    \begin{lstlisting}[language=Python]
print(shopping_list)
# Output: ['milk', 'eggs', 'whole wheat bread']
    \end{lstlisting}
    
\end{frame}

\begin{frame}[fragile]
    \frametitle{Modifying List Elements: Important Rules}
    
    \begin{enumerate}
        \item \textbf{Modification is In-Place}: The change happens directly to the original list. No new list is created.
        
        \pause
        
        \item \textbf{Index Must Exist}: You can only assign a value to an index that is already in the list.
    \end{enumerate}
    
    \begin{alertblock}{IndexError}
        The following code will cause an \texttt{IndexError} because index 3 does not exist in a 3-element list (valid indices are 0, 1, 2).
\begin{lstlisting}[language=Python]
# THIS WILL CAUSE AN ERROR!
shopping_list[3] = 'butter'
\end{lstlisting}
    \end{alertblock}
    
    \pause
    
    \begin{block}{Negative Indices Work Too!}
        You can also use negative indices to modify elements from the end of the list.
\begin{lstlisting}[language=Python]
# Another way to change the last item
shopping_list[-1] = 'sourdough bread'
\end{lstlisting}
    \end{block}
    
\end{frame}

\begin{frame}[fragile]
    \frametitle{Getting a Subset: Slicing (1/3) - The Concept}
    
    Slicing is a powerful technique to get a sub-list, or a 'slice', from an original list.
    
    \begin{itemize}
        \item It allows you to access a range of elements at once.
        \item \textbf{Important:} Slicing creates a \textit{new list}; the original list is not modified.
    \end{itemize}
    
    \bigskip
    
    \textbf{Syntax}: \texttt{list\_name[start:stop]}
    
    \begin{alertblock}{The "Golden Rule" of Slicing}
        The slice includes the element at the \texttt{start} index up to, but \textbf{not including}, the element at the \texttt{stop} index.
    \end{alertblock}
    
\end{frame}

\begin{frame}[fragile]
    \frametitle{Getting a Subset: Slicing (2/3) - A Visual Example}
    
    Let's see how the slice \texttt{letters[1:4]} works.
    
    \begin{block}{Our List and its Indices}
\begin{lstlisting}[language=Python]
# List:
letters = ['a', 'b', 'c', 'd', 'e']
# Index:    0    1    2    3    4 
\end{lstlisting}
    \end{block}
    
    \begin{exampleblock}{Applying the Rule: \texttt{letters[1:4]}}
        \begin{enumerate}
            \item \textbf{Start} at index \texttt{1} (element \texttt{'b'}).
            \item \textbf{Continue} through indices \texttt{2} and \texttt{3}.
            \item \textbf{Stop} \textit{before} reaching index \texttt{4}.
        \end{enumerate}
    \end{exampleblock}
    
    \textbf{Result}: A new list is returned.
\begin{lstlisting}[language=Python]
letters[1:4]  # Returns ['b', 'c', 'd']
\end{lstlisting}
    
\end{frame}

\begin{frame}[fragile]
    \frametitle{Getting a Subset: Slicing (3/3) - Shortcuts}
    
    Python provides handy shortcuts by letting you omit the \texttt{start} or \texttt{stop} index.
    
    \begin{itemize}
        \item \textbf{From the beginning:} \texttt{[:stop]}
        \begin{itemize}
            \item Omitting \texttt{start} defaults to the beginning (index 0).
            \item \texttt{letters[:3]} \textit{\# Returns ['a', 'b', 'c']}
        \end{itemize}
        
        \bigskip
        
        \item \textbf{To the end:} \texttt{[start:]}
        \begin{itemize}
            \item Omitting \texttt{stop} defaults to the end of the list.
            \item \texttt{letters[2:]} \textit{\# Returns ['c', 'd', 'e']}
        \end{itemize}
        
        \bigskip
        
        \item \textbf{Copy the whole list:} \texttt{[:]}
        \begin{itemize}
            \item A common idiom to create a \textbf{shallow copy} of a list.
            \item \texttt{letters\_copy = letters[:]}
        \end{itemize}
    \end{itemize}
    
\end{frame}

\begin{frame}[fragile]
    \frametitle{Adding Elements to a List (1/2): Append}
    
    One of the most powerful features of Python lists is that they are \textbf{mutable}, meaning you can change them after they are created. Let's see how to grow a list.

    \begin{block}{1. Append to the End: \texttt{.append()}}
        The \texttt{.append()} method adds a single item to the very end of the list. It modifies the list in-place.
        \begin{itemize}
            \item \textbf{Purpose}: Add one item to the end.
            \item \textbf{Syntax}: \texttt{list\_name.append(item\_to\_add)}
        \end{itemize}
    \end{block}

    \pause

    \begin{exampleblock}{Code Example}
\begin{lstlisting}[language=Python]
# We start with a simple to-do list
tasks = ['laundry']

# Add 'dishes' to the end of the list
tasks.append('dishes')

# tasks is now ['laundry', 'dishes']
\end{lstlisting}
    \end{exampleblock}

\end{frame}

\begin{frame}[fragile]
    \frametitle{Adding Elements to a List (2/2): Insert}
    
    What if you need to add an item to the beginning or middle of a list? For precise control over placement, we use the \texttt{.insert()} method.

    \begin{block}{2. Insert at a Specific Index: \texttt{.insert()}}
        The \texttt{.insert()} method adds an item at a specific index, shifting all subsequent elements to the right.
        \begin{itemize}
            \item \textbf{Purpose}: Add one item at a specific position.
            \item \textbf{Syntax}: \texttt{list\_name.insert(index, item\_to\_add)}
        \end{itemize}
    \end{block}

    \pause

    \begin{exampleblock}{Code Example}
\begin{lstlisting}[language=Python]
# Continuing with our list from before:
# tasks is ['laundry', 'dishes']

# An urgent task comes up, add it to the beginning (index 0)
tasks.insert(0, 'email boss')

# tasks is now ['email boss', 'laundry', 'dishes']
\end{lstlisting}
    \end{exampleblock}

\end{frame}

\begin{frame}[fragile]
    \frametitle{Removing Elements from a List (1/3): The \texttt{del} Statement}
    
    There are several ways to remove items, depending on what you know (the index or the value).
    
    \begin{block}{1. The \texttt{del} Statement}
        Removes an item at a known index. The item is permanently deleted and cannot be accessed.
        
\begin{lstlisting}[language=Python]
# Initial list
tasks = ['email boss', 'laundry', 'dishes', 'call mom']

# Remove the item at index 1
del tasks[1]

# tasks is now: ['email boss', 'dishes', 'call mom']
\end{lstlisting}
    \end{block}
    
    \begin{itemize}
        \item Use \texttt{del} when you want to "delete and forget" an item at a specific position.
    \end{itemize}
    
\end{frame}

\begin{frame}[fragile]
    \frametitle{Removing Elements from a List (2/3): Methods}
    
    \begin{block}{2. The \texttt{.pop()} Method}
        Removes the item at a specific index and lets you \textbf{use it}. If no index is given, it removes the \textbf{last} item.
\begin{lstlisting}[language=Python]
# tasks = ['email boss', 'laundry', 'dishes', 'call mom']
done_task = tasks.pop(0)

# done_task is now: 'email boss'
# tasks is now: ['laundry', 'dishes', 'call mom']
\end{lstlisting}
    \end{block}
    
    \pause
    
    \begin{block}{3. The \texttt{.remove()} Method}
        Removes the \textbf{first occurrence} of a specific value. A \texttt{ValueError} occurs if the item is not found.
\begin{lstlisting}[language=Python]
# tasks = ['email boss', 'laundry', 'dishes', 'call mom']
tasks.remove('dishes')

# tasks is now: ['email boss', 'laundry', 'call mom']
\end{lstlisting}
    \end{block}
    
\end{frame}

\begin{frame}[fragile]
    \frametitle{Removing Elements from a List (3/3): Summary}
    
    \begin{table}
        \centering
        \caption{Which removal method should you use?}
        \begin{tabular}{l l p{6cm}}
            \hline
            \textbf{Goal} & \textbf{Syntax} & \textbf{Why?} \\
            \hline \hline
            Remove by index (discard) & \texttt{del my\_list[i]} & Simple and direct for removal by index. \\
            \hline
            Remove by index (and use) & \texttt{item = my\_list.pop(i)} & Returns the removed item for you to use. \\
            \hline
            Remove last item (and use) & \texttt{item = my\_list.pop()} & A common, useful default for the \texttt{.pop()} method. \\
            \hline
            Remove by value & \texttt{my\_list.remove('x')} & The only method that searches for an item by its value. \\
            \hline
        \end{tabular}
    \end{table}
    
\end{frame}

\begin{frame}[fragile]
    \frametitle{Common List Functions and Operators - The len() Function}
    
    The \texttt{len()} function is a fundamental tool that returns the number of elements in a list. It simply inspects the list without modifying it.
    
    \begin{itemize}
        \item \textbf{Purpose}: To determine the size or length of a list.
        \item \textbf{Analogy}: Counting the number of items in your shopping cart.
        \item \textbf{Syntax}: \texttt{len(your\_list)}
    \end{itemize}
    
    \begin{block}{Example: Counting Tasks}
\begin{lstlisting}[language=Python]
# Example: Counting tasks
tasks = ['wash car', 'buy groceries', 'pay bills']
num_tasks = len(tasks)
print(f"You have {num_tasks} tasks remaining.")

# You can also check if a list is empty
if len(tasks) == 0:
    print("All done!")
\end{lstlisting}
    \vspace{2mm}
    \textbf{Output:}
\begin{verbatim}
You have 3 tasks remaining.
\end{verbatim}
    \end{block}
    
\end{frame}

\begin{frame}[fragile]
    \frametitle{Common List Functions and Operators - Operators}

    \begin{block}{1. Concatenation (+): Combining Lists}
        The \texttt{+} operator joins two lists end-to-end to create a \textbf{new} list.
        \begin{itemize}
            \item \textbf{Key Point}: This operation is \textbf{non-destructive}; the original lists are not changed.
        \end{itemize}
\begin{lstlisting}[language=Python]
morning_tasks = ['gym', 'check email']
evening_tasks = ['cook dinner', 'read book']

all_day_tasks = morning_tasks + evening_tasks
print(f"All tasks: {all_day_tasks}")
print(f"Morning tasks are still: {morning_tasks}")
\end{lstlisting}
    \textbf{Output:}
\begin{verbatim}
All tasks: ['gym', 'check email', 'cook dinner', 'read book']
Morning tasks are still: ['gym', 'check email']
\end{verbatim}
    \end{block}

    \pause

    \begin{block}{2. Repetition (*): Repeating a List}
        The \texttt{*} operator creates a \textbf{new} list by repeating its elements a specified number of times.
        \begin{itemize}
            \item \textbf{Use Case}: Excellent for initializing a list with a default value.
        \end{itemize}
\begin{lstlisting}[language=Python]
# Example 1: A simple cheer
cheer = ['Go!'] * 3
print(cheer)

# Example 2: Initializing a game board row
row = ['_'] * 3
print(f"Empty row: {row}")
\end{lstlisting}
    \textbf{Output:}
\begin{verbatim}
['Go!', 'Go!', 'Go!']
Empty row: ['_', '_', '_']
\end{verbatim}
    \end{block}
\end{frame}

\begin{frame}[fragile]
    \frametitle{Chapter Summary: Mastering Lists}
    
    This chapter gave us the tools to store, access, and manage collections of data.
    
    \begin{block}{1. The Anatomy of a List}
        A container for holding multiple items in a specific sequence.
        \begin{itemize}
            \item \textbf{Ordered}: The position of each item is preserved.
            \item \textbf{Mutable}: The list can be changed after it's created.
            \item \textbf{Syntax}: Defined using square brackets \texttt{[]}.
        \end{itemize}
    \end{block}
    
    \pause
    
    \begin{block}{2. Accessing \& Modifying Data}
        \begin{itemize}
            \item \textbf{Indexing}: Get or change a \textit{single} element. Remember it starts at zero! \\ \texttt{my\_list[0] = 'new\_value'}
            \item \textbf{Slicing}: Get a \textit{new sub-list} containing a range of elements. \\ \texttt{sub\_list = my\_list[0:2]}
        \end{itemize}
    \end{block}
    
    \pause
    
    \begin{block}{3. Managing Content with Methods}
        \begin{itemize}
            \item \textbf{Adding Elements}: \texttt{.append(item)}, \texttt{.insert(index, item)}
            \item \textbf{Removing Elements}: \texttt{.pop(index)}, \texttt{.remove(value)}
        \end{itemize}
    \end{block}
    
\end{frame}

\begin{frame}[fragile]
    \frametitle{Next Steps: Putting Lists to Work}
    
    \begin{block}{The Challenge}
        So far, we've interacted with lists one element at a time.
        \vspace{0.5em}
        
        How do we perform an action on \textit{every single item} in a list of 1,000 participants without writing 1,000 lines of code?
    \end{block}
    
    \pause
    
    \begin{alertblock}{The Solution: Loops}
        In the next chapter, we will introduce \textbf{loops}. By combining lists and loops, you'll be able to:
        \begin{itemize}
            \item Iterate through every item in a list automatically.
            \item Perform calculations on entire datasets.
            \item Search for items that meet specific criteria.
        \end{itemize}
    \end{alertblock}
    
    \vspace{1em}
    
    \centering
    \textit{This combination is the key to unlocking data processing and automation in Python.}
    
\end{frame}


\end{document}